% ============================================================
% 架空論文Ｘ　都市部の若年層における自動車購入前の情報探索と販売店訪問
%   基礎演習Ⅱ 第4回「文献講読と批判的思考」の AI 演習用。
%   記入例（worksheets/answers/04.tex）が参照している頁に合わせてある：
%     問い p.2 ／ 方法（購入経験者に限定）p.4 ／ 表3（情報源別）p.7 ／
%     主な結果 p.8 ／ 関心度別の分析 p.9 ／ 限界 p.12
%   ビルド：make paper-x.pdf（Makefile を参照）
% ============================================================
% ============================================================
% 基礎演習Ⅱ　授業用の架空論文　共通プリアンブル
%   学術誌の体裁をまねた1段組の論文レイアウト。各論文（paper-x.tex など）の
%   冒頭で % ============================================================
% 基礎演習Ⅱ　授業用の架空論文　共通プリアンブル
%   学術誌の体裁をまねた1段組の論文レイアウト。各論文（paper-x.tex など）の
%   冒頭で % ============================================================
% 基礎演習Ⅱ　授業用の架空論文　共通プリアンブル
%   学術誌の体裁をまねた1段組の論文レイアウト。各論文（paper-x.tex など）の
%   冒頭で \input{preamble} して使う。LuaLaTeX（jlreq）で組む。
% ============================================================
\documentclass[paper=a4paper,fontsize=10pt,line_length=42zw,number_of_lines=38,
  baselineskip=1.7zw,head_space=22mm,foot_space=24mm]{jlreq}

\usepackage{xcolor}
\usepackage{tabularray}
\UseTblrLibrary{booktabs}
\usepackage{caption}
\usepackage{enumitem}
\usepackage{needspace}
\usepackage{fancyhdr}
\usepackage{pgfplots}
\pgfplotsset{compat=1.18}
\usepackage[hidelinks]{hyperref}

\definecolor{ink}{HTML}{1D2330}
\definecolor{inktwo}{HTML}{4A5261}
\definecolor{inkthree}{HTML}{757D8A}
\definecolor{rulegray}{HTML}{B9BEB4}
\definecolor{papertwo}{HTML}{ECEEE8}
\definecolor{bad}{HTML}{B3261E}
\definecolor{badtint}{HTML}{F9DEDA}
\definecolor{barone}{HTML}{2747C9}
\definecolor{bartwo}{HTML}{D68A17}

% ------------------------------------------------------------
% 論文ごとに設定する情報（各論文の冒頭で \renewcommand する）
% ------------------------------------------------------------
\newcommand*\paperlabel{架空論文}         % 架空論文Ｘ など
\newcommand*\journalname{（架空の学術誌）}
\newcommand*\journalvol{}
\newcommand*\runninghead{}        % 柱に出す著者名など

% ------------------------------------------------------------
% ページスタイル
% ------------------------------------------------------------
\pagestyle{fancy}
\fancyhf{}
\renewcommand{\headrulewidth}{0.3pt}
\renewcommand{\footrulewidth}{0pt}
\fancyhead[L]{\small\color{inktwo}\journalname　\journalvol}
\fancyhead[R]{\small\color{inktwo}\runninghead}
\fancyfoot[R]{\small\thepage}
\fancyfoot[L]{\scriptsize\color{inkthree}\sffamily\paperlabel　基礎演習Ⅱ（名古屋市立大学 経済学部）の教材として作成した架空の論文。著者・調査・データ・文献はすべて架空。}
\fancypagestyle{firstpage}{%
  \fancyhf{}
  \renewcommand{\headrulewidth}{0pt}
  \fancyfoot[R]{\small\thepage}
  \fancyfoot[L]{\scriptsize\color{inkthree}\sffamily\paperlabel　基礎演習Ⅱ（名古屋市立大学 経済学部）の教材として作成した架空の論文。著者・調査・データ・文献はすべて架空。}}

% ------------------------------------------------------------
% 見出し・キャプション・箇条書き
% ------------------------------------------------------------
\ModifyHeading{section}{font=\normalsize\sffamily\bfseries,before_space=1.4\baselineskip,after_space=0.5\baselineskip,
  lines=2}
\ModifyHeading{subsection}{font=\normalsize\sffamily\bfseries,before_space=0.9\baselineskip,after_space=0.3\baselineskip,
  lines=1}
\ModifyHeading{subsubsection}{font=\normalsize\sffamily,before_space=0.6\baselineskip,after_space=0.2\baselineskip,
  lines=1}
\renewcommand{\thesection}{\arabic{section}.}
\renewcommand{\thesubsection}{\arabic{section}.\arabic{subsection}}
\renewcommand{\thesubsubsection}{\arabic{section}.\arabic{subsection}.\arabic{subsubsection}}

\DeclareCaptionLabelFormat{jp}{#1#2}
\captionsetup{labelformat=jp,labelsep=quad,font=small,labelfont={sf,bf},justification=centering,skip=4pt}
\setlist{nosep,leftmargin=1.8em}

% ------------------------------------------------------------
% 表題部　\makepapertitle{題名}{副題}{著者}{所属}
% ------------------------------------------------------------
\newcommand\makepapertitle[4]{%
  \thispagestyle{firstpage}
  \begingroup
  \noindent{\small\color{inktwo}\journalname　\journalvol}\hfill{\small\sffamily\color{inktwo}【\paperlabel】}\par
  \vspace{0.5mm}{\color{inktwo}\rule{\linewidth}{0.4pt}}\par
  \vspace{7mm}
  \begin{center}
    {\Large\sffamily\bfseries #1}\par\vspace{1.5mm}
    {\normalsize\sffamily #2}\par\vspace{6mm}
    {\normalsize #3}\par\vspace{1mm}
    {\small\color{inktwo}#4}\par
  \end{center}
  \vspace{4mm}
  \endgroup}

% 要旨とキーワード
\newenvironment{abstractbox}{%
  \begin{list}{}{\leftmargin=2\zw\rightmargin=2\zw}\item\relax\small
  {\sffamily\bfseries 要旨}\par\vspace{0.5mm}}{\end{list}\vspace{2mm}}
\newcommand\keywords[1]{{\small\hspace{2\zw}\textbf{キーワード}：#1}\par\vspace{2mm}}

% 架空であることの注記（第1ページの表題の下に置く）
\newcommand\fictionnote[1]{%
  \noindent\colorbox{badtint}{\parbox{\dimexpr\linewidth-2\fboxsep\relax}{\footnotesize\color{bad}\sffamily #1}}\par\vspace{3mm}}

% 表の共通スタイル（booktabs 風）
\NewTblrEnviron{ptab}
\SetTblrInner[ptab]{cells={font=\footnotesize},rowsep=1.2pt,colsep=5pt}
\newcommand\tnote[1]{\par\vspace{1mm}\noindent\begin{minipage}{\linewidth}\footnotesize\color{inktwo}注：#1\end{minipage}\par}

% 本文の流れの中に置く表と図（フロートにしない。頁の位置を固定するため）
%   \begin{inlinetable}{キャプション} ... \end{inlinetable}
%   \begin{inlinefigure}{キャプション} ... \end{inlinefigure}（キャプションは下）
\newenvironment{inlinetable}[1]{%
  \par\vspace{0.6\baselineskip}\noindent\begin{minipage}{\linewidth}\captionof{table}{#1}\centering\footnotesize}%
  {\end{minipage}\par\vspace{0.8\baselineskip}}
\newenvironment{inlinefigure}[1]{%
  \par\vspace{0.6\baselineskip}\noindent\begin{minipage}{\linewidth}\centering\def\figcap{#1}}%
  {\captionof{figure}{\figcap}\end{minipage}\par\vspace{0.8\baselineskip}}

% 有意水準の記号
\newcommand\sig[1]{\textsuperscript{#1}}

% 引用文献リスト
\newenvironment{references}{\begin{list}{}{\leftmargin=2\zw\itemindent=-2\zw\itemsep=1.5pt\parsep=0pt}\small}{\end{list}}
 して使う。LuaLaTeX（jlreq）で組む。
% ============================================================
\documentclass[paper=a4paper,fontsize=10pt,line_length=42zw,number_of_lines=38,
  baselineskip=1.7zw,head_space=22mm,foot_space=24mm]{jlreq}

\usepackage{xcolor}
\usepackage{tabularray}
\UseTblrLibrary{booktabs}
\usepackage{caption}
\usepackage{enumitem}
\usepackage{needspace}
\usepackage{fancyhdr}
\usepackage{pgfplots}
\pgfplotsset{compat=1.18}
\usepackage[hidelinks]{hyperref}

\definecolor{ink}{HTML}{1D2330}
\definecolor{inktwo}{HTML}{4A5261}
\definecolor{inkthree}{HTML}{757D8A}
\definecolor{rulegray}{HTML}{B9BEB4}
\definecolor{papertwo}{HTML}{ECEEE8}
\definecolor{bad}{HTML}{B3261E}
\definecolor{badtint}{HTML}{F9DEDA}
\definecolor{barone}{HTML}{2747C9}
\definecolor{bartwo}{HTML}{D68A17}

% ------------------------------------------------------------
% 論文ごとに設定する情報（各論文の冒頭で \renewcommand する）
% ------------------------------------------------------------
\newcommand*\paperlabel{架空論文}         % 架空論文Ｘ など
\newcommand*\journalname{（架空の学術誌）}
\newcommand*\journalvol{}
\newcommand*\runninghead{}        % 柱に出す著者名など

% ------------------------------------------------------------
% ページスタイル
% ------------------------------------------------------------
\pagestyle{fancy}
\fancyhf{}
\renewcommand{\headrulewidth}{0.3pt}
\renewcommand{\footrulewidth}{0pt}
\fancyhead[L]{\small\color{inktwo}\journalname　\journalvol}
\fancyhead[R]{\small\color{inktwo}\runninghead}
\fancyfoot[R]{\small\thepage}
\fancyfoot[L]{\scriptsize\color{inkthree}\sffamily\paperlabel　基礎演習Ⅱ（名古屋市立大学 経済学部）の教材として作成した架空の論文。著者・調査・データ・文献はすべて架空。}
\fancypagestyle{firstpage}{%
  \fancyhf{}
  \renewcommand{\headrulewidth}{0pt}
  \fancyfoot[R]{\small\thepage}
  \fancyfoot[L]{\scriptsize\color{inkthree}\sffamily\paperlabel　基礎演習Ⅱ（名古屋市立大学 経済学部）の教材として作成した架空の論文。著者・調査・データ・文献はすべて架空。}}

% ------------------------------------------------------------
% 見出し・キャプション・箇条書き
% ------------------------------------------------------------
\ModifyHeading{section}{font=\normalsize\sffamily\bfseries,before_space=1.4\baselineskip,after_space=0.5\baselineskip,
  lines=2}
\ModifyHeading{subsection}{font=\normalsize\sffamily\bfseries,before_space=0.9\baselineskip,after_space=0.3\baselineskip,
  lines=1}
\ModifyHeading{subsubsection}{font=\normalsize\sffamily,before_space=0.6\baselineskip,after_space=0.2\baselineskip,
  lines=1}
\renewcommand{\thesection}{\arabic{section}.}
\renewcommand{\thesubsection}{\arabic{section}.\arabic{subsection}}
\renewcommand{\thesubsubsection}{\arabic{section}.\arabic{subsection}.\arabic{subsubsection}}

\DeclareCaptionLabelFormat{jp}{#1#2}
\captionsetup{labelformat=jp,labelsep=quad,font=small,labelfont={sf,bf},justification=centering,skip=4pt}
\setlist{nosep,leftmargin=1.8em}

% ------------------------------------------------------------
% 表題部　\makepapertitle{題名}{副題}{著者}{所属}
% ------------------------------------------------------------
\newcommand\makepapertitle[4]{%
  \thispagestyle{firstpage}
  \begingroup
  \noindent{\small\color{inktwo}\journalname　\journalvol}\hfill{\small\sffamily\color{inktwo}【\paperlabel】}\par
  \vspace{0.5mm}{\color{inktwo}\rule{\linewidth}{0.4pt}}\par
  \vspace{7mm}
  \begin{center}
    {\Large\sffamily\bfseries #1}\par\vspace{1.5mm}
    {\normalsize\sffamily #2}\par\vspace{6mm}
    {\normalsize #3}\par\vspace{1mm}
    {\small\color{inktwo}#4}\par
  \end{center}
  \vspace{4mm}
  \endgroup}

% 要旨とキーワード
\newenvironment{abstractbox}{%
  \begin{list}{}{\leftmargin=2\zw\rightmargin=2\zw}\item\relax\small
  {\sffamily\bfseries 要旨}\par\vspace{0.5mm}}{\end{list}\vspace{2mm}}
\newcommand\keywords[1]{{\small\hspace{2\zw}\textbf{キーワード}：#1}\par\vspace{2mm}}

% 架空であることの注記（第1ページの表題の下に置く）
\newcommand\fictionnote[1]{%
  \noindent\colorbox{badtint}{\parbox{\dimexpr\linewidth-2\fboxsep\relax}{\footnotesize\color{bad}\sffamily #1}}\par\vspace{3mm}}

% 表の共通スタイル（booktabs 風）
\NewTblrEnviron{ptab}
\SetTblrInner[ptab]{cells={font=\footnotesize},rowsep=1.2pt,colsep=5pt}
\newcommand\tnote[1]{\par\vspace{1mm}\noindent\begin{minipage}{\linewidth}\footnotesize\color{inktwo}注：#1\end{minipage}\par}

% 本文の流れの中に置く表と図（フロートにしない。頁の位置を固定するため）
%   \begin{inlinetable}{キャプション} ... \end{inlinetable}
%   \begin{inlinefigure}{キャプション} ... \end{inlinefigure}（キャプションは下）
\newenvironment{inlinetable}[1]{%
  \par\vspace{0.6\baselineskip}\noindent\begin{minipage}{\linewidth}\captionof{table}{#1}\centering\footnotesize}%
  {\end{minipage}\par\vspace{0.8\baselineskip}}
\newenvironment{inlinefigure}[1]{%
  \par\vspace{0.6\baselineskip}\noindent\begin{minipage}{\linewidth}\centering\def\figcap{#1}}%
  {\captionof{figure}{\figcap}\end{minipage}\par\vspace{0.8\baselineskip}}

% 有意水準の記号
\newcommand\sig[1]{\textsuperscript{#1}}

% 引用文献リスト
\newenvironment{references}{\begin{list}{}{\leftmargin=2\zw\itemindent=-2\zw\itemsep=1.5pt\parsep=0pt}\small}{\end{list}}
 して使う。LuaLaTeX（jlreq）で組む。
% ============================================================
\documentclass[paper=a4paper,fontsize=10pt,line_length=42zw,number_of_lines=38,
  baselineskip=1.7zw,head_space=22mm,foot_space=24mm]{jlreq}

\usepackage{xcolor}
\usepackage{tabularray}
\UseTblrLibrary{booktabs}
\usepackage{caption}
\usepackage{enumitem}
\usepackage{needspace}
\usepackage{fancyhdr}
\usepackage{pgfplots}
\pgfplotsset{compat=1.18}
\usepackage[hidelinks]{hyperref}

\definecolor{ink}{HTML}{1D2330}
\definecolor{inktwo}{HTML}{4A5261}
\definecolor{inkthree}{HTML}{757D8A}
\definecolor{rulegray}{HTML}{B9BEB4}
\definecolor{papertwo}{HTML}{ECEEE8}
\definecolor{bad}{HTML}{B3261E}
\definecolor{badtint}{HTML}{F9DEDA}
\definecolor{barone}{HTML}{2747C9}
\definecolor{bartwo}{HTML}{D68A17}

% ------------------------------------------------------------
% 論文ごとに設定する情報（各論文の冒頭で \renewcommand する）
% ------------------------------------------------------------
\newcommand*\paperlabel{架空論文}         % 架空論文Ｘ など
\newcommand*\journalname{（架空の学術誌）}
\newcommand*\journalvol{}
\newcommand*\runninghead{}        % 柱に出す著者名など

% ------------------------------------------------------------
% ページスタイル
% ------------------------------------------------------------
\pagestyle{fancy}
\fancyhf{}
\renewcommand{\headrulewidth}{0.3pt}
\renewcommand{\footrulewidth}{0pt}
\fancyhead[L]{\small\color{inktwo}\journalname　\journalvol}
\fancyhead[R]{\small\color{inktwo}\runninghead}
\fancyfoot[R]{\small\thepage}
\fancyfoot[L]{\scriptsize\color{inkthree}\sffamily\paperlabel　基礎演習Ⅱ（名古屋市立大学 経済学部）の教材として作成した架空の論文。著者・調査・データ・文献はすべて架空。}
\fancypagestyle{firstpage}{%
  \fancyhf{}
  \renewcommand{\headrulewidth}{0pt}
  \fancyfoot[R]{\small\thepage}
  \fancyfoot[L]{\scriptsize\color{inkthree}\sffamily\paperlabel　基礎演習Ⅱ（名古屋市立大学 経済学部）の教材として作成した架空の論文。著者・調査・データ・文献はすべて架空。}}

% ------------------------------------------------------------
% 見出し・キャプション・箇条書き
% ------------------------------------------------------------
\ModifyHeading{section}{font=\normalsize\sffamily\bfseries,before_space=1.4\baselineskip,after_space=0.5\baselineskip,
  lines=2}
\ModifyHeading{subsection}{font=\normalsize\sffamily\bfseries,before_space=0.9\baselineskip,after_space=0.3\baselineskip,
  lines=1}
\ModifyHeading{subsubsection}{font=\normalsize\sffamily,before_space=0.6\baselineskip,after_space=0.2\baselineskip,
  lines=1}
\renewcommand{\thesection}{\arabic{section}.}
\renewcommand{\thesubsection}{\arabic{section}.\arabic{subsection}}
\renewcommand{\thesubsubsection}{\arabic{section}.\arabic{subsection}.\arabic{subsubsection}}

\DeclareCaptionLabelFormat{jp}{#1#2}
\captionsetup{labelformat=jp,labelsep=quad,font=small,labelfont={sf,bf},justification=centering,skip=4pt}
\setlist{nosep,leftmargin=1.8em}

% ------------------------------------------------------------
% 表題部　\makepapertitle{題名}{副題}{著者}{所属}
% ------------------------------------------------------------
\newcommand\makepapertitle[4]{%
  \thispagestyle{firstpage}
  \begingroup
  \noindent{\small\color{inktwo}\journalname　\journalvol}\hfill{\small\sffamily\color{inktwo}【\paperlabel】}\par
  \vspace{0.5mm}{\color{inktwo}\rule{\linewidth}{0.4pt}}\par
  \vspace{7mm}
  \begin{center}
    {\Large\sffamily\bfseries #1}\par\vspace{1.5mm}
    {\normalsize\sffamily #2}\par\vspace{6mm}
    {\normalsize #3}\par\vspace{1mm}
    {\small\color{inktwo}#4}\par
  \end{center}
  \vspace{4mm}
  \endgroup}

% 要旨とキーワード
\newenvironment{abstractbox}{%
  \begin{list}{}{\leftmargin=2\zw\rightmargin=2\zw}\item\relax\small
  {\sffamily\bfseries 要旨}\par\vspace{0.5mm}}{\end{list}\vspace{2mm}}
\newcommand\keywords[1]{{\small\hspace{2\zw}\textbf{キーワード}：#1}\par\vspace{2mm}}

% 架空であることの注記（第1ページの表題の下に置く）
\newcommand\fictionnote[1]{%
  \noindent\colorbox{badtint}{\parbox{\dimexpr\linewidth-2\fboxsep\relax}{\footnotesize\color{bad}\sffamily #1}}\par\vspace{3mm}}

% 表の共通スタイル（booktabs 風）
\NewTblrEnviron{ptab}
\SetTblrInner[ptab]{cells={font=\footnotesize},rowsep=1.2pt,colsep=5pt}
\newcommand\tnote[1]{\par\vspace{1mm}\noindent\begin{minipage}{\linewidth}\footnotesize\color{inktwo}注：#1\end{minipage}\par}

% 本文の流れの中に置く表と図（フロートにしない。頁の位置を固定するため）
%   \begin{inlinetable}{キャプション} ... \end{inlinetable}
%   \begin{inlinefigure}{キャプション} ... \end{inlinefigure}（キャプションは下）
\newenvironment{inlinetable}[1]{%
  \par\vspace{0.6\baselineskip}\noindent\begin{minipage}{\linewidth}\captionof{table}{#1}\centering\footnotesize}%
  {\end{minipage}\par\vspace{0.8\baselineskip}}
\newenvironment{inlinefigure}[1]{%
  \par\vspace{0.6\baselineskip}\noindent\begin{minipage}{\linewidth}\centering\def\figcap{#1}}%
  {\captionof{figure}{\figcap}\end{minipage}\par\vspace{0.8\baselineskip}}

% 有意水準の記号
\newcommand\sig[1]{\textsuperscript{#1}}

% 引用文献リスト
\newenvironment{references}{\begin{list}{}{\leftmargin=2\zw\itemindent=-2\zw\itemsep=1.5pt\parsep=0pt}\small}{\end{list}}

\renewcommand*\paperlabel{架空論文Ｘ}
\renewcommand*\journalname{消費と流通の研究}
\renewcommand*\journalvol{第18巻 第2号（2026年）}
\renewcommand*\runninghead{鳴海・柏木：自動車購入前の情報探索と販売店訪問}

\begin{document}

\makepapertitle{都市部の若年層における自動車購入前の情報探索と販売店訪問}
{――オンライン情報探索の量は訪問店舗数と商談時間をどう変えるか――}
{鳴海　千尋\textsuperscript{1}　・　柏木　遼\textsuperscript{2}}
{\textsuperscript{1}東海商科大学 商学部　\textsuperscript{2}都市消費研究所}

\fictionnote{この論文は、名古屋市立大学 経済学部「基礎演習Ⅱ」第4回「文献講読と批判的思考」の教材として作成した架空の論文です。著者、所属、掲載誌、調査、データ、引用文献はすべて架空であり、実在の研究・団体・人物とは関係ありません。数値は授業の練習のために作ったもので、実際の市場の状況を示すものではありません。}

\begin{abstractbox}
若年層の「車離れ」が語られる一方で、都市部の若年層の多くは、自動車を購入する前に動画サイトや比較サイトで詳しく情報を集めるようになった。本研究は、購入前のオンライン情報探索の量が、販売店への訪問行動とどのように関係しているかを明らかにする。東京23区および政令指定都市に居住し、過去3年以内に自動車を購入した20〜34歳の1,042人を対象にWeb調査を行い、情報探索量と、訪問した販売店の数、延べ訪問回数、来店1回あたりの商談時間との関係を分析した。その結果、情報探索量が多い回答者ほど訪問店舗数と延べ訪問回数は少ないが、来店1回あたりの商談時間は長いことが示された。情報源別にみると、動画サイトの利用が多い回答者ほど訪問店舗数が少なかった。これらの関係は、自動車への関心度の高低にかかわらず認められた。以上の結果は、若年層にとっての販売店が、複数の候補を「比べる場」から、絞り込んだ候補を確かめて「決める場」へと変わりつつあることを示唆する。
\end{abstractbox}
\keywords{情報探索、販売店訪問、若年層、ウェブルーミング、自動車購買行動}

% ------------------------------------------------------------
\section{はじめに}

\subsection{背景}

日本では2000年代後半から、若年層の「車離れ」がくり返し語られてきた。運転免許の取得率や自動車保有率の低下、都市部での公共交通への依存、所得の伸び悩みなどが、その背景として挙げられる（仮屋, 2019; 東原・三宅, 2021）。しかし、都市部に住む若年層のすべてが自動車を買わなくなったわけではない。結婚や出産、転居、趣味といった生活の節目で自動車を購入する若年層は一定数存在し、その購入のしかたは、上の世代とは大きく異なっている。

変化の中心にあるのは、購入前の情報環境である。かつて自動車の購入を検討する人は、カタログを取り寄せ、複数の販売店を回って実車を見比べ、営業担当者の説明を聞きながら候補を絞っていった。今日では、メーカーや販売店の公式サイトに加え、価格比較サイト、口コミサイト、SNS、そして試乗やレビューの動画が、購入前の情報の大半を提供している。とりわけ動画サイトには、車種ごとの内装の質感や走行の様子、維持費の試算まで、かつては販売店に行かなければ得られなかった情報が公開されている（瀬川, 2024）。

こうした情報環境の変化は、販売店の役割に直接かかわる。自動車販売の実務家からは、「来店が減った」「来店したときにはもう車種が決まっている」という声が聞かれるようになった（全国自動車小売研究会, 2024）。しかし、この実感が、購入前のオンライン情報探索の量とどのように関係しているのかを、都市部の若年層について定量的に示した研究は少ない。来店の減少が「販売店が不要になった」ことを意味するのか、それとも「販売店で行うことの内容が変わった」ことを意味するのかによって、販売店のとるべき対応は大きく異なる。本研究はこの点を、購入者本人への調査から明らかにしようとするものである。

\subsection{本研究の問い}

本研究の問いは、次の二つである。

\begin{description}[style=nextline,leftmargin=2\zw,labelsep=0.5\zw,font=\normalfont\bfseries,itemsep=2pt]
  \item[RQ1] 都市部の若年層において、購入前のオンライン情報探索の量は、購入までに訪問した販売店の数および延べ訪問回数と、どのように関係しているか。
  \item[RQ2] 購入前のオンライン情報探索の量は、来店1回あたりの商談時間と、どのように関係しているか。
\end{description}

これらに加えて、補助的な問いとして、利用した情報源の種類（公式サイト、比較サイト、口コミ、動画、SNSなど）によって、販売店訪問との関係が異なるかどうかを検討する（RQ3）。

本研究の焦点は、情報探索の「量」と販売店訪問の「量」（店舗数、回数、時間）の関係にある。来店して何をしたのか、つまり来店の目的や内容については、本研究の調査では直接には尋ねていない。この点は、本研究の範囲の限界として第6節で改めて述べる。

\subsection{本稿の構成}

第2節で先行研究を整理し、仮説を導く。第3節で調査と測定の方法を、第4節で結果を示す。第5節で結果を解釈し、販売店への含意を論じる。第6節で本研究の限界と今後の課題を述べる。

% ------------------------------------------------------------
\section{先行研究と仮説}

\subsection{購入前の情報探索}

消費者が購入前に行う情報探索は、記憶の中から知識を取り出す内部探索と、外部の情報源にあたる外部探索に分けられる（大崎, 2008）。外部探索の量は、探索によって得られる便益（よりよい選択ができること）と、探索にかかる費用（時間、労力、金銭）のバランスで決まると考えられてきた。自動車は高額で購入頻度が低く、選択を誤ったときの損失が大きいため、外部探索の量が多い製品の代表とされる。

インターネットの普及は、外部探索の費用を大きく引き下げた。従来、自動車の外部探索の多くは、カタログの入手と販売店の訪問という形で行われ、訪問そのものが探索費用の大きな部分を占めていた。オンラインで情報が得られるようになると、探索費用のうち「移動と対面」にかかる部分が削られ、探索の総量は増える一方で、探索の「場所」が販売店からオンラインへ移ると予想される（桐生, 2020）。このとき、販売店の訪問は探索の手段としての意味を弱め、別の役割を担うようになる可能性がある。

\subsection{オンライン情報探索と店舗訪問}

小売業の研究では、店舗で実物を確かめてからオンラインで購入する「ショールーミング」と、オンラインで調べてから店舗で購入する「ウェブルーミング」が対比されてきた（久保田, 2018）。自動車は、試乗や現車の確認、契約の手続き、登録や納車といった対面の要素が多く、オンラインだけで購入が完結することはまれである。したがって自動車の購買は典型的なウェブルーミング型の過程であり、オンラインでの探索が店舗訪問を「置き換える」のではなく、店舗訪問の「前に」置かれる構造をもつ。

ウェブルーミング型の購買では、オンラインでの事前の絞り込みが進むほど、店舗で比較する候補が減り、訪問する店舗の数が減ることが、家電や住宅の研究で報告されている（森下・安藤, 2022; Hartley \& Nkemelu, 2019）。自動車についても、販売店への訪問回数が長期的に減少しているという業界統計がある（全国自動車小売研究会, 2024）。しかし、これらの研究と統計は、訪問の「数」の減少を示す一方で、訪問1回あたりの内容の変化については扱っていない。

一方、絞り込みが進んだ状態での来店は、比較のための来店とは性質が異なると考えられる。候補を絞り込んだ消費者は、来店時に試乗、見積、下取りの査定、契約条件の確認など、購入を決めるための複数の行為を1回の来店でまとめて行おうとする。この場合、来店1回あたりの滞在時間は長くなる可能性がある。この点は米国の販売店に関する実務調査で示唆されているが（Greenfield Automotive Insights, 2023）、日本の若年層について学術的に検討した研究は見当たらない。

\subsection{仮説}

以上から、本研究では次の仮説を設定する。

\begin{description}[style=nextline,leftmargin=2\zw,labelsep=0.5\zw,font=\normalfont\bfseries,itemsep=2pt]
  \item[H1] 購入前のオンライン情報探索の量が多いほど、購入までに訪問した販売店の数は少ない。
  \item[H2] 購入前のオンライン情報探索の量が多いほど、購入までの延べ訪問回数は少ない。
  \item[H3] 購入前のオンライン情報探索の量が多いほど、来店1回あたりの商談時間は長い。
\end{description}

情報源の種類に関しては、実車の外観や内装、走行の様子を視覚的に確認できる動画サイトは、「見るための来店」を代替する度合いが大きいと考えられる（Lindqvist, 2021）。そこで探索的に次の仮説を置く。

\begin{description}[style=nextline,leftmargin=2\zw,labelsep=0.5\zw,font=\normalfont\bfseries,itemsep=2pt]
  \item[H4] 動画サイトの利用が多いほど、購入までに訪問した販売店の数は少ない。
\end{description}

なお、情報探索の量と店舗訪問の関係を解釈するうえでは、自動車への関心度（製品関与）が両者に共通する要因となっている可能性に注意が必要である（相原, 2023）。関心度の高い消費者は情報探索を多く行う一方で、販売店を訪れること自体を楽しむため、店舗数も多くなるかもしれない。この場合、関心度を考慮しなければ、探索量と店舗数の負の関係は見かけ上弱まるか、隠される。本研究では、関心度を統制変数として投入するとともに、関心度の高低で分けた分析を行う。

% ------------------------------------------------------------
\clearpage
\section{方法}

\subsection{調査の概要}

本研究のデータは、2025年11月に実施したWeb調査による。調査は、調査会社のインターネット調査モニター（登録者約180万人）を利用し、次の三つの条件をすべて満たす人を対象とした。

\begin{enumerate}[label=(\arabic*),leftmargin=2.5\zw,itemsep=1pt]
  \item 東京23区または政令指定都市（20市）に居住していること。本研究では、これらの地域をまとめて「都市部」と呼ぶ。
  \item 調査時点で20〜34歳であること。
  \item 過去3年以内（2022年11月〜2025年10月）に新車または中古車の乗用車を購入し、その購入の主な決定者であったこと。
\end{enumerate}

条件(3)により、本研究の対象は、自動車を実際に購入した人に限られる。購入を検討したが購入に至らなかった人や、購入を検討していない人は含まれない。これは、販売店への訪問行動を購入過程の始まりから終わりまで回顧的に尋ねるためであり、購入に至らなかった人については検討の起点と終点を定めにくいと判断した。したがって、本研究の結果は購入者における情報探索と販売店訪問の関係を示すものであり、都市部の若年層全体の来店行動の変化を示すものではない。

スクリーニング調査には14,206人が回答し、このうち条件(1)〜(3)をすべて満たした1,318人に本調査を依頼した。1,112人が本調査に回答し、回答時間が極端に短いもの（中央値の3分の1未満）、同じ選択肢を連続して選んだもの、訪問店舗数や商談時間に非現実的な値を記入したもの（訪問店舗数が20店以上、商談時間が1回あたり10時間以上）を除外した。最終的な有効回答は1,042人である。年齢と性別の割り付けは行っていない。

\subsection{測定}

\subsubsection*{(1) オンライン情報探索量}

購入を検討し始めてから契約するまでの期間について、次の二つを尋ねた。第一に、利用した情報源の種類である。表3に示す8種類のオンライン情報源を提示し、該当するものをすべて選んでもらった（0〜8種類）。第二に、探索時間である。「購入を検討し始めてから契約するまでに、インターネットで自動車の情報を調べた時間は、合計でどのくらいでしたか」と尋ね、「1時間未満」から「50時間以上」までの8段階で回答を得て、各段階の中央値を時間に換算した。両者は中程度に相関していた（$r = .46$）ため、それぞれを標準化して平均した値を「情報探索量得点」とし、主な分析に用いた。また、群間の比較のために、得点の3分位で回答者を低群、中群、高群に分けた。

\subsubsection*{(2) 販売店訪問}

同じ期間について、(a) 訪問した販売店の数（異なる店舗の数。同じ店舗に複数回行った場合は1と数える）、(b) 延べ訪問回数、(c) 来店1回あたりの平均的な滞在時間（分）を尋ねた。(c) は「販売店に1回行ったとき、店にいた時間は平均してどのくらいでしたか」という質問で測定し、本稿では「商談時間」と呼ぶ。この時間には、商談のほか試乗や待ち時間が含まれうる。

\subsubsection*{(3) 自動車への関心度}

「あなたは自動車にどのくらい関心がありますか」と尋ね、「まったく関心がない」から「非常に関心がある」までの5件法で回答を得た。関心度別の分析では、4または5と回答した人を高関心群、1〜3と回答した人を低関心群とした。

\subsubsection*{(4) 統制変数}

性別、年齢、購入した車が新車か中古車か、購入価格帯（5区分）、初めての自動車購入かどうか、運転免許取得からの年数、世帯年収（6区分）を尋ねた。調査票の主な項目は付録に示す。

\subsection{分析の方法}

まず、情報探索量の3群別に販売店訪問の3指標（訪問店舗数、延べ訪問回数、商談時間）の平均を比較し、一元配置分散分析で群差を検定した。次に、情報源の種類ごとに利用者と非利用者の訪問店舗数を比較した。さらに、3指標を従属変数、情報探索量得点と統制変数を独立変数とする重回帰分析を行い、最後に関心度の高低で回答者を分けて同じ分析を群ごとに行った。世帯年収を「回答しない」とした32人は、重回帰分析から除いた。

% ------------------------------------------------------------
\section{結果}

\subsection{回答者の属性}

表1に回答者の属性を示す。男性が58.2\%、年齢は30〜34歳が40.4\%で最も多く、購入した車は新車が54.3\%であった。初めての自動車購入であった人が61.0\%と過半数を占め、購入価格帯は100〜300万円未満に58.4\%が集中している。自動車への関心度（5件法）は、高関心群（4または5）が487人（46.7\%）、低関心群（1〜3）が555人（53.3\%）であった。

\begin{inlinetable}{回答者の属性（$n = 1{,}042$）}
\begin{minipage}[t]{0.47\linewidth}
\begin{ptab}{colspec={Q[l] Q[l] Q[r] Q[r]},row{1}={font=\footnotesize\sffamily\bfseries}}
\toprule
項目 & 区分 & 人数 & \% \\
\midrule
性別 & 男性 & 606 & 58.2 \\
 & 女性 & 436 & 41.8 \\
年齢 & 20〜24歳 & 230 & 22.1 \\
 & 25〜29歳 & 391 & 37.5 \\
 & 30〜34歳 & 421 & 40.4 \\
購入した車 & 新車 & 566 & 54.3 \\
 & 中古車 & 476 & 45.7 \\
初めての購入 & はい & 636 & 61.0 \\
 & いいえ & 406 & 39.0 \\
免許取得後 & 3年未満 & 171 & 16.4 \\
の年数 & 3〜7年未満 & 365 & 35.0 \\
 & 7年以上 & 506 & 48.6 \\
\bottomrule
\end{ptab}
\end{minipage}\hfill
\begin{minipage}[t]{0.50\linewidth}
\begin{ptab}{colspec={Q[l] Q[l] Q[r] Q[r]},row{1}={font=\footnotesize\sffamily\bfseries}}
\toprule
項目 & 区分 & 人数 & \% \\
\midrule
購入価格帯 & 100万円未満 & 198 & 19.0 \\
 & 100〜200万円未満 & 322 & 30.9 \\
 & 200〜300万円未満 & 287 & 27.5 \\
 & 300〜400万円未満 & 148 & 14.2 \\
 & 400万円以上 & 87 & 8.4 \\
世帯年収 & 300万円未満 & 146 & 14.0 \\
 & 300〜500万円未満 & 312 & 29.9 \\
 & 500〜700万円未満 & 283 & 27.2 \\
 & 700〜1,000万円未満 & 198 & 19.0 \\
 & 1,000万円以上 & 71 & 6.8 \\
 & 回答しない & 32 & 3.1 \\
\bottomrule
\end{ptab}
\end{minipage}
\end{inlinetable}

\subsection{情報探索量と販売店訪問}

回答者全体では、利用した情報源は平均3.9種類、探索時間は平均16.8時間（中央値12時間）であった。販売店訪問については、訪問店舗数の平均が2.1店、延べ訪問回数の平均が3.4回、商談時間の平均が62.3分であった。

表2は、情報探索量の3群別に、販売店訪問の3指標を比較したものである。訪問店舗数は、低群の2.6店から高群の1.6店へと、探索量が多い群ほど少なかった（$F(2, 1039) = 62.4$, $p < .001$）。延べ訪問回数も、低群の3.9回から高群の3.0回へと少なくなっていた（$F(2, 1039) = 21.7$, $p < .001$）。一方、商談時間は、低群の48.3分から高群の77.6分へと、探索量が多い群ほど長かった（$F(2, 1039) = 98.5$, $p < .001$）。

\begin{inlinetable}{情報探索量の3群別にみた販売店訪問}
\begin{ptab}{colspec={Q[l] Q[c] Q[c] Q[c] Q[c] Q[c]},row{1}={font=\footnotesize\sffamily\bfseries},row{2}={font=\footnotesize\sffamily\bfseries}}
\toprule
 & 低群 & 中群 & 高群 & 全体 & $F$ 値 \\
 & ($n = 347$) & ($n = 348$) & ($n = 347$) & ($n = 1{,}042$) & \\
\midrule
情報源の種類数（種類） & 2.1 & 3.9 & 5.7 & 3.9 & --- \\
探索時間（時間） & 4.6 & 13.2 & 32.5 & 16.8 & --- \\
\midrule
訪問店舗数（店） & 2.6 (1.4) & 2.1 (1.1) & 1.6 (0.9) & 2.1 (1.2) & 62.4\sig{***} \\
延べ訪問回数（回） & 3.9 (2.2) & 3.4 (1.9) & 3.0 (1.7) & 3.4 (2.0) & 21.7\sig{***} \\
商談時間（分） & 48.3 (22.5) & 61.0 (25.8) & 77.6 (31.2) & 62.3 (29.4) & 98.5\sig{***} \\
\midrule
1店舗あたりの訪問回数（回） & 1.5 & 1.6 & 1.9 & 1.6 & --- \\
店にいた時間の合計（分） & 188 & 207 & 233 & 212 & --- \\
\bottomrule
\end{ptab}
\tnote{括弧内は標準偏差。情報源の種類数と探索時間は、群分けに用いた変数の平均。1店舗あたりの訪問回数は延べ訪問回数を訪問店舗数で割った値、店にいた時間の合計は延べ訪問回数と商談時間の積（いずれも各群の平均から算出）。\sig{***} $p < .001$。}
\end{inlinetable}

表2の下段には、これらから導かれる二つの指標を示した。1店舗あたりの訪問回数（延べ訪問回数÷訪問店舗数）は、低群の1.5回から高群の1.9回へと増えている。つまり、探索量の多い回答者は、少ない店舗にくり返し通っている。また、店にいた時間の合計（延べ訪問回数×商談時間）は、低群の188分に対して高群は233分と、むしろ高群のほうが長い。探索量の多い回答者は、販売店で過ごす時間が短いのではなく、その時間を少数の店舗に集中させているといえる。図1は、訪問店舗数と商談時間の群差を図示したものである。

\begin{inlinefigure}{情報探索量の3群別にみた訪問店舗数（左）と商談時間（右）}
\begin{tikzpicture}
\begin{axis}[
  width=60mm,height=46mm,ybar,bar width=7mm,
  ymin=0,ymax=3.2,ylabel={訪問店舗数（店）},
  symbolic x coords={低群,中群,高群},xtick=data,
  enlarge x limits=0.25,
  nodes near coords={\pgfmathprintnumber[fixed,fixed zerofill,precision=1]{\pgfplotspointmeta}},nodes near coords style={font=\footnotesize},
  tick label style={font=\footnotesize},label style={font=\footnotesize},
  ymajorgrids,grid style={rulegray!60},axis line style={rulegray},
  every axis plot/.append style={fill=barone,draw=none}]
\addplot coordinates {(低群,2.6) (中群,2.1) (高群,1.6)};
\end{axis}
\end{tikzpicture}\hspace{6mm}
\begin{tikzpicture}
\begin{axis}[
  width=60mm,height=46mm,ybar,bar width=7mm,
  ymin=0,ymax=95,ylabel={商談時間（分）},
  symbolic x coords={低群,中群,高群},xtick=data,
  enlarge x limits=0.25,
  nodes near coords={\pgfmathprintnumber[fixed,fixed zerofill,precision=1]{\pgfplotspointmeta}},nodes near coords style={font=\footnotesize},
  tick label style={font=\footnotesize},label style={font=\footnotesize},
  ymajorgrids,grid style={rulegray!60},axis line style={rulegray},
  every axis plot/.append style={fill=bartwo,draw=none}]
\addplot coordinates {(低群,48.3) (中群,61.0) (高群,77.6)};
\end{axis}
\end{tikzpicture}
\end{inlinefigure}

\subsection{情報源の種類別の分析}

表3は、8種類のオンライン情報源について、利用率と、利用者・非利用者それぞれの訪問店舗数の平均を示したものである。利用率が最も高いのはメーカーの公式サイト（82.4\%）で、販売店の公式サイト・SNS（61.7\%）、価格比較・中古車情報サイト（58.9\%）が続く。動画サイトの利用率は46.3\%で、回答者のほぼ半数が購入前に動画を見ていた。

利用者と非利用者の訪問店舗数の差が最も大きかったのは動画サイトであり、利用者の1.7店に対して非利用者は2.4店であった（$p < .001$）。動画サイトの利用頻度別にみると、利用なしの回答者の訪問店舗数が2.4店、月に数回以下の利用者が1.9店、週1回以上の利用者が1.5店と、利用が多いほど訪問店舗数は少なかった（傾向性の検定 $p < .001$）。H4は支持された。

これに対して、価格比較・中古車情報サイトは、利用者のほうが非利用者よりも訪問店舗数が多い唯一の情報源であった（2.2店対2.0店, $p = .012$）。メーカーの公式サイトと自動車専門メディアでは、利用者と非利用者の差は有意でなかった。

利用した情報源の種類数でみると、0〜2種類の回答者の訪問店舗数が2.7店、3〜4種類が2.1店、5種類以上が1.6店と、種類数が多いほど訪問店舗数は少なかった。動画サイトを除いた7種類で同じ集計をしても、この傾向は変わらなかった（2.6店、2.1店、1.7店）。訪問店舗数の少なさは、動画サイトの利用だけで説明されるものではなく、情報源を幅広く参照することと結びついている。

なお、情報源は単独ではなく組み合わせて使われている。動画サイトの利用者の91.3\%はメーカーの公式サイトを、68.5\%は口コミ・レビューサイトを併用していた。動画サイトだけを利用した回答者は11人にとどまり、動画は他の情報源で得た候補を確かめる段階で使われていると考えられる。

\begin{inlinetable}{情報源の種類別にみた利用率と訪問店舗数}
\begin{ptab}{colspec={Q[l] Q[r] Q[r] Q[c] Q[c] Q[c] Q[c]},row{1}={font=\footnotesize\sffamily\bfseries},row{2}={font=\footnotesize\sffamily\bfseries}}
\toprule
情報源 & 利用者数 & 利用率 & \SetCell[c=2]{c} 訪問店舗数（店） & & 差 & $p$ \\
 & （人） & （\%） & 利用者 & 非利用者 & & \\
\midrule
メーカーの公式サイト & 859 & 82.4 & 2.1 & 2.2 & $-0.1$ & .21 \\
販売店の公式サイト・SNS & 643 & 61.7 & 2.0 & 2.3 & $-0.3$ & $<.001$ \\
価格比較・中古車情報サイト & 614 & 58.9 & 2.2 & 2.0 & $+0.2$ & .012 \\
口コミ・レビューサイト & 517 & 49.6 & 2.0 & 2.2 & $-0.2$ & .008 \\
動画サイト（試乗・レビュー動画） & 482 & 46.3 & 1.7 & 2.4 & $-0.7$ & $<.001$ \\
SNS（X、Instagram など） & 397 & 38.1 & 1.9 & 2.2 & $-0.3$ & $<.001$ \\
自動車専門のニュースメディア & 287 & 27.5 & 2.0 & 2.1 & $-0.1$ & .34 \\
知人からのオンラインでの紹介 & 259 & 24.9 & 1.8 & 2.2 & $-0.4$ & $<.001$ \\
\bottomrule
\end{ptab}
\tnote{$n = 1{,}042$。複数回答。差は利用者の平均から非利用者の平均を引いた値。$p$ 値は利用者と非利用者の平均の差の $t$ 検定による。}
\end{inlinetable}

\newpage
\subsection{重回帰分析}

表4は、販売店訪問の3指標を従属変数とする重回帰分析の結果である。統制変数を投入したうえでも、情報探索量得点は訪問店舗数に対して負（$\beta = -.31$, $p < .001$）、延べ訪問回数に対して負（$\beta = -.18$, $p < .001$）、商談時間に対して正（$\beta = .27$, $p < .001$）の有意な係数を示した。すなわち、\textbf{情報探索量が多い回答者ほど訪問店舗数と延べ訪問回数は少ないが、来店1回あたりの商談時間は長い}。H1、H2、H3はいずれも支持された。3指標の中では、訪問店舗数に対する探索量の影響が最も大きく、延べ訪問回数に対する影響は相対的に小さかった。これは、探索量の多い回答者が少数の店舗にくり返し通っているという表2の結果と整合する。

統制変数についてみると、自動車への関心度は、訪問店舗数、延べ訪問回数、商談時間のいずれにも正の係数を示した。関心度の高い回答者は、より多くの店舗に、より多くの回数行き、1回の来店時間も長い。初めての自動車購入である回答者も、3指標すべてで正の係数を示した。新車の購入者は、中古車の購入者に比べて訪問店舗数が少なく、商談時間が長かった。

係数の実質的な大きさをみると、情報探索量得点が1標準偏差高い回答者は、他の条件が同じなら、訪問店舗数が0.37店少なく、延べ訪問回数が0.36回少なく、商談時間が7.9分長い。訪問店舗数の平均が2.1店であることを考えると、探索量による差は無視できない大きさである。

なお、関心度を投入しないモデルでは、情報探索量得点の訪問店舗数に対する係数は $\beta = -.33$ であり、関心度を投入した表4のモデル（$-.31$）との差はわずかであった。独立変数間の分散拡大係数（VIF）は最大で1.6であり、多重共線性の問題は認められなかった。

\begin{inlinetable}{販売店訪問の3指標を従属変数とする重回帰分析（標準化係数）}
\begin{ptab}{colspec={Q[l] Q[c] Q[c] Q[c]},row{1}={font=\footnotesize\sffamily\bfseries}}
\toprule
独立変数 & 訪問店舗数 & 延べ訪問回数 & 商談時間 \\
\midrule
情報探索量得点 & $-.31$\sig{***} & $-.18$\sig{***} & $.27$\sig{***} \\
自動車への関心度 & $.12$\sig{**} & $.15$\sig{***} & $.09$\sig{*} \\
女性（ダミー） & $-.03$ & $.01$ & $.04$ \\
年齢 & $-.05$ & $-.04$ & $.02$ \\
新車（ダミー） & $-.08$\sig{*} & $.06$ & $.11$\sig{**} \\
購入価格帯 & $.07$\sig{*} & $.09$\sig{*} & $.10$\sig{**} \\
初めての自動車購入（ダミー） & $.10$\sig{**} & $.12$\sig{**} & $.14$\sig{***} \\
運転免許取得からの年数 & $-.02$ & $-.03$ & $-.01$ \\
世帯年収 & $.02$ & $.03$ & $-.01$ \\
\midrule
調整済み $R^2$ & .17 & .11 & .15 \\
$n$ & 1,010 & 1,010 & 1,010 \\
\bottomrule
\end{ptab}
\tnote{世帯年収を「回答しない」とした32人を除く。\sig{*} $p < .05$、\sig{**} $p < .01$、\sig{***} $p < .001$。}
\end{inlinetable}

\newpage
\subsection{自動車への関心度別の分析}

第2節で述べたように、自動車への関心度が探索量と店舗訪問の両方を高める共通の要因であるなら、探索量と訪問店舗数の負の関係は、関心度の違いによって生じた見かけの関係である可能性がある。この可能性を確かめるため、関心度の高低で回答者を分け、群ごとに重回帰分析を行った（表5）。

表5の上段に示すように、高関心群は低関心群に比べて、情報探索量得点が高く、訪問店舗数も多く、商談時間も長い。関心度の高い回答者が、情報も多く集め、店舗も多く回るという想定は、平均の水準としては当てはまる。しかし、群内の関係をみると、探索量得点は高関心群でも低関心群でも、訪問店舗数に対して負（それぞれ $\beta = -.26$, $-.34$）、商談時間に対して正（$\beta = .24$, $.29$）の有意な係数を示した。全体のモデルに探索量得点と関心度の交互作用項を加えても、係数は有意でなかった（訪問店舗数 $\beta = .04$、商談時間 $\beta = -.03$、いずれも n.s.）。すなわち、探索量が多いほど訪問店舗数が少なく商談時間が長いという関係は、関心度の高低にかかわらず認められ、関心度の違いだけで説明することはできない。

\begin{inlinetable}{自動車への関心度別にみた情報探索量と販売店訪問}
\begin{ptab}{colspec={Q[l] Q[c] Q[c]},row{1}={font=\footnotesize\sffamily\bfseries},row{2}={font=\footnotesize\sffamily\bfseries}}
\toprule
 & 高関心群 & 低関心群 \\
 & ($n = 471$) & ($n = 539$) \\
\midrule
情報探索量得点（平均） & $0.21$ & $-0.18$ \\
訪問店舗数（店、平均） & 2.3 & 1.9 \\
延べ訪問回数（回、平均） & 3.7 & 3.1 \\
商談時間（分、平均） & 66.8 & 58.4 \\
\midrule
探索量得点 → 訪問店舗数（$\beta$） & $-.26$\sig{***} & $-.34$\sig{***} \\
探索量得点 → 延べ訪問回数（$\beta$） & $-.14$\sig{**} & $-.21$\sig{***} \\
探索量得点 → 商談時間（$\beta$） & $.24$\sig{***} & $.29$\sig{***} \\
\bottomrule
\end{ptab}
\tnote{関心度は5件法の1項目で測定し、4・5を高関心群、1〜3を低関心群とした。下段は表4と同じ統制変数（関心度を除く）を投入した群別の重回帰分析における情報探索量得点の標準化係数。世帯年収の無回答者を除く。\sig{**} $p < .01$、\sig{***} $p < .001$。}
\end{inlinetable}

\subsection{補足的な分析}

結果の頑健性を確かめるため、三つの補足的な分析を行った。第一に、新車の購入者（$n = 549$）と中古車の購入者（$n = 461$）に分けて表4と同じ重回帰分析を行ったところ、情報探索量得点の訪問店舗数に対する係数はそれぞれ $\beta = -.29$、$-.33$、商談時間に対する係数は $.25$、$.28$ であり、いずれも有意であった（$p < .001$）。第二に、初めての自動車購入である回答者とそうでない回答者に分けた場合も、係数の符号と有意性は変わらなかった。第三に、情報探索量得点を構成する二つの変数（情報源の種類数と探索時間）を、合成せずに別々に投入した。訪問店舗数に対する係数は、種類数が $\beta = -.22$、探索時間が $-.15$ であり、どちらも有意であったが、種類数の係数のほうが大きかった。商談時間に対しても同様に、種類数（$.19$）が探索時間（$.13$）を上回った。販売店訪問との関係は、オンラインで費やした時間の長さよりも、参照した情報源の幅広さとより強く結びついているといえる。

% ------------------------------------------------------------
\section{考察}

\subsection{結果の要約}

本研究の結果は、次の三点にまとめられる。第一に、都市部の若年層の購入者において、購入前のオンライン情報探索の量が多いほど、訪問した販売店の数と延べ訪問回数は少ない（H1、H2）。第二に、探索量が多いほど、来店1回あたりの商談時間は長く、店にいた時間の合計はむしろ長い（H3）。第三に、情報源の中では動画サイトの利用が訪問店舗数の少なさと最も強く関係し（H4）、価格比較サイトの利用はむしろ訪問店舗数の多さと関係していた。これらの関係は、自動車への関心度の高低にかかわらず認められた。

\subsection{「比べる場」から「決める場」へ}

これらの結果は、若年層にとっての販売店の意味が変わりつつあることを示唆する。第2節で述べたように、オンライン化は探索費用のうち「移動と対面」にかかる部分を引き下げた。複数の車種を見比べ、価格や装備を比較し、評判を確かめるという、かつては販売店を回って行っていた作業の多くは、オンラインで済ませられるようになった。その結果、販売店を訪れる時点では候補がすでに絞られており、来店は複数の候補を「比べる場」ではなく、絞り込んだ候補を確かめて「決める場」になっていると解釈できる。

この解釈は、探索量の多い回答者が少数の店舗にくり返し通い、1回あたりの商談時間が長いという結果と整合する。候補を絞り込んだ消費者は、1回の来店で試乗、見積、下取りの査定、契約条件の確認といった複数の行為をまとめて行おうとするため、滞在時間が長くなると考えられる。ただし、本研究では来店時に何をしたのかを直接には尋ねていないため、商談時間が長い理由についての説明は推測にとどまる。来店の目的や内容を測定し、「決める場」という解釈を直接に検証することは、今後の課題である。

探索費用の観点からみると、訪問店舗数の減少は探索の「減少」ではなく、探索の「配分の変化」として理解できる。探索量の多い回答者は、情報源の種類も探索時間も多く、販売店で過ごす時間の合計も長い。探索は全体として減ったのではなく、比較の段階がオンラインへ、確認と決定の段階が販売店へと、役割ごとに配分されるようになったと考えられる。

\subsection{情報源による違い}

動画サイトの利用が訪問店舗数の少なさと最も強く関係していたことは、動画が「見るための来店」を代替しているという仮説（H4）と整合する。内装の質感、荷室の広さ、走行の様子といった視覚的な情報は、文字情報では伝わりにくく、かつては実車を見るために販売店を訪れる主な理由であった。動画はこの部分を代替することで、来店の必要性を減らしていると考えられる。

一方、価格比較サイトの利用者がむしろ多くの店舗を訪れていたことは、情報源の性質によってオンライン探索と店舗訪問の関係が異なることを示している。価格比較サイトは候補を絞り込むというよりも、複数の販売店から見積を取り、価格を交渉する行動を促している可能性がある。この場合、オンラインの情報は来店を減らすのではなく、来店先を選ぶための情報として用いられている。情報源の種類ごとに、来店との関係の方向が異なることは、オンライン情報探索を一括りに扱う議論の限界を示すものである。

\subsection{販売店への含意}

本研究の結果は、販売店にとって二つの含意をもつ。第一に、若年層の来店数の減少は、来店の価値の低下を意味しない。探索量の多い若年層の購入者は、店舗数を絞ったうえで、1回の来店により長い時間を使っている。販売店は来店数の減少を嘆くよりも、少なくなった来店1回あたりの質を高めることに注力すべきである。具体的には、来店時に候補がすでに絞られていることを前提に、初回の来店で試乗から見積、下取りの査定、契約条件の提示までを一度に進められる体制が求められる。

第二に、来店前の情報探索の段階で販売店自身が情報を提供することの重要性である。動画サイトでの探索が訪問店舗数を減らしているのであれば、販売店が自店の在庫や試乗の様子を動画で公開することは、絞り込みの段階で候補に残るための手段になる。販売店の公式サイト・SNSの利用者が非利用者よりも訪問店舗数が少なかったことは、販売店の発信が絞り込みに影響している可能性を示唆する。

一方で、価格比較サイトの利用者が複数の店舗を訪れていたことは、価格の比較が依然として来店の動機になっていることを示す。来店前に価格の目安が共有されている以上、販売店は見積の内容を来店前に提示し、来店時には価格以外の価値（試乗の体験、納車後の対応、下取りの条件など）を伝えることが求められる。来店の回数が減る中で、来店前の情報提供と来店時の対応は、切り離して考えることができない。

\subsection{理論的含意}

探索費用の低下が外部探索の総量を増やすという従来の予測に加えて、本研究は、探索の「場所」と「段階」の配分が変わることを示した。オンライン探索と店舗訪問は、代替か補完かという二者択一ではなく、購買過程の段階ごとに異なる役割を担う。この見方は、ウェブルーミング型の購買を、オンラインから店舗への単なる移行としてではなく、段階ごとの役割分担として捉え直すことを求める。自動車のように対面の要素が多い製品では、この役割分担が特に明確に現れると考えられる。

また、探索量と店舗訪問の関係が自動車への関心度の高低にかかわらず認められたことは、この役割分担が、製品関与の高い一部の消費者に限られた現象ではないことを示している。関心度の高い消費者は全体として多く探索し多く来店するが、その中でも探索の多い人ほど来店先を絞っている。探索の「場所」の配分は、関与の水準とは独立に、情報環境の変化によって規定されていると考えられる。

% ------------------------------------------------------------
\newpage
\section{本研究の限界と今後の課題}

本研究にはいくつかの限界がある。

第一に、データがWeb調査モニターから得られたものである点である。インターネット調査モニターは、インターネットの利用時間が長く、オンラインでの情報収集に慣れた人に偏っている可能性がある。この偏りは、本研究の中心的な変数であるオンライン情報探索量を過大に見積もる方向に働きうる。都市部の若年層の購入者の全体像を示すためには、調査方法の異なるデータによる確認が必要である。

第二に、本研究は一時点の横断調査であり、探索量と販売店訪問の関係は相関関係にとどまる。「情報探索を多く行ったから訪問店舗数が少なくなった」のか、「もともと少数の店舗で決めるつもりだった人が、その分オンラインで多く調べた」のかを、本研究のデータでは区別できない。また、購入過程を終えた後に回顧的に尋ねているため、探索時間や商談時間の回答には記憶の誤りが含まれうる。因果関係を明らかにするには、購入の検討中の人を追跡する縦断的な調査が必要である。

第三に、第3節で述べたように、本研究の対象は自動車を実際に購入した人に限られる。購入を検討したが購入しなかった人の販売店訪問は、本研究の範囲の外にある。したがって、本研究の結果から販売店全体の来店の変化を論じることはできない。

第四に、本研究は販売店訪問の「量」を測定したが、来店の目的や内容は測定していない。商談時間が長い理由として本稿が述べた解釈（1回の来店で試乗から契約までを進める）は推測であり、来店時の行動を直接に尋ねる調査によって検証する必要がある。

第五に、本研究は東京23区と政令指定都市を一括して「都市部」として扱い、地域による違いを検討していない。公共交通の整備の状況や自動車への依存度は、同じ都市部でも地域によって大きく異なる。自動車の利用が多い地域では、販売店訪問のあり方が本研究の結果と異なる可能性がある。地域別の分析と、都市部以外の地域との比較は、今後の課題である。

% ------------------------------------------------------------
\section{おわりに}

本研究は、都市部の若年層の自動車購入者を対象に、購入前のオンライン情報探索の量と販売店訪問の関係を検討した。情報探索量が多い購入者ほど訪問する販売店の数は少ないが、来店1回あたりの商談時間は長く、この関係は自動車への関心度の高低にかかわらず認められた。若年層にとっての販売店は、候補を「比べる場」から、絞り込んだ候補を確かめて「決める場」へと変わりつつある。来店の目的や内容、地域による違い、購入に至らなかった人の行動を含めた検討が、今後の課題である。

% ------------------------------------------------------------
\newpage
\section*{引用文献}

{\footnotesize\color{bad}\sffamily 注意：以下の引用文献はすべて架空であり、実在しない。}

\begin{references}
  \item 相原亮（2023）「製品関与と情報探索の再検討――高関与製品のオンライン探索を中心に」『マーケティング科学論集』31, 55--72。
  \item 大崎玲（2008）『消費者の情報探索行動』都市経済出版。
  \item 仮屋修一（2019）「若年層の自動車保有の決定要因――三大都市圏の比較」『交通と社会』32(1), 45--61。
  \item 桐生麻衣（2020）「オンライン化による探索費用の変化と購買前行動」『消費と流通の研究』12(2), 23--40。
  \item 久保田慎吾（2018）「ショールーミングとウェブルーミングの規定要因」『商業学論集』54(3), 1--18。
  \item 瀬川ひかり（2024）「動画メディアと若年層の購買前行動」『情報行動研究』7(1), 12--29。
  \item 全国自動車小売研究会（2024）『販売店来店動向レポート2024』全国自動車小売研究会。
  \item 東原和也・三宅奈々（2021）「都市部若年層の自動車利用と購入意向」『都市生活研究』9, 77--95。
  \item 森下恵・安藤隆（2022）「住宅購入における事前の情報探索と見学件数」『住宅市場研究』15(1), 101--118。
  \item Greenfield Automotive Insights (2023). \textit{Dealer Visit Patterns among Young Buyers: 2023 Report}. Greenfield Automotive Insights.
  \item Hartley, J., \& Nkemelu, A. (2019). Prepurchase screening and store visits in consumer durables. \textit{Journal of Retail Channel Studies}, 27(4), 311--329.
  \item Lindqvist, M. (2021). Video reviews as substitutes for physical inspection: Evidence from durable goods. \textit{International Review of Consumer Research}, 14(2), 88--104.
\end{references}

% ------------------------------------------------------------
\section*{付録　調査票の主な項目}

\begin{description}[style=nextline,leftmargin=1\zw,font=\normalfont\sffamily\bfseries,itemsep=3pt]
  \item[スクリーニング]
  Q1 お住まいの市区町村（選択）。Q2 年齢（数値）。Q3 2022年11月以降に、新車または中古車の乗用車を購入しましたか（はい／いいえ）。Q4 その購入について、あなたは主な決定者でしたか（はい／いいえ／共同で決めた）。※「共同で決めた」は対象に含めた。
  \item[情報探索]
  Q5 購入を検討し始めてから契約するまでに、インターネットで利用した情報源をすべて選んでください（メーカーの公式サイト／販売店の公式サイト・SNS／価格比較・中古車情報サイト／口コミ・レビューサイト／動画サイト（試乗・レビュー動画）／SNS（X、Instagram など）／自動車専門のニュースメディア／知人からのオンラインでの紹介（メッセージアプリなど）／いずれも利用していない）。
  Q6 購入を検討し始めてから契約するまでに、インターネットで自動車の情報を調べた時間は、合計でどのくらいでしたか（1時間未満／1〜3時間未満／3〜5時間未満／5〜10時間未満／10〜20時間未満／20〜30時間未満／30〜50時間未満／50時間以上）。
  Q7 動画サイトで自動車の情報を見た頻度（利用していない／月に数回以下／週に1回以上）。
  \item[販売店訪問]
  Q8 購入を検討し始めてから契約するまでに、訪問した販売店はいくつありましたか。同じ店に何度か行った場合は1と数えてください（数値）。
  Q9 販売店に行った回数の合計（数値）。
  Q10 販売店に1回行ったとき、店にいた時間は平均してどのくらいでしたか（分、数値）。
  \item[関心度]
  Q11 あなたは自動車にどのくらい関心がありますか（1 まったく関心がない／2 あまり関心がない／3 どちらともいえない／4 やや関心がある／5 非常に関心がある）。
  \item[属性]
  性別、購入した車（新車／中古車）、購入価格帯（5区分）、初めての自動車購入か（はい／いいえ）、運転免許取得からの年数、世帯年収（6区分、「回答しない」を含む）。
\end{description}

\vspace{6mm}
\begin{center}
{\footnotesize\color{inkthree}（受理：2026年3月／この論文は架空であり、受理日も架空である）}
\end{center}

\end{document}
